% !TeX program = xelatex
% 电工杯数学建模竞赛模板 — 基于 neepumcm 文档类
\documentclass[a4paper]{neepumcm}
% ⛔ graphicx, booktabs, colortbl, xcolor 已由 neepumcm.cls 加载，不要重复加载
\usepackage{tikz}
\usetikzlibrary{arrows.meta,positioning,shapes.geometric,fit,calc}
\usepackage{algorithm2e}
\usepackage{float}
\usepackage[section]{placeins}
\usepackage{needspace}
\usepackage{longtable}

% === 引用上标格式 ===
\makeatletter
\def\@cite#1#2{\textsuperscript{[{#1\if@tempswa , #2\fi}]}}
\makeatother

\mcmsetup{CTeX = true,
        tcn = [报名序号],
        problem = [A/B],
        timu = [论文题目]}
\setmainfont{Times New Roman}
\setmonofont[
    Path=fonts/,
    UprightFont = *-Regular,
    BoldFont = *-Bold,
    ItalicFont = *-Italic
]{UbuntuMono}

\begin{document}

\begin{abstract}

[中文摘要内容：问题概述 + 每个子问题的方法和数值结果 + 结论]

\begin{keywords}
[关键词1] \quad [关键词2] \quad [关键词3]
\end{keywords}
\end{abstract}

\maketitle
\newpage
\tableofcontents
\clearpage

% === 正文 ===
\input{sections/1_restatement}
\input{sections/2_analysis}
\input{sections/3_assumptions}
\input{sections/4_symbols}
\input{sections/5_problem1}
\input{sections/6_problem2}
\input{sections/7_problem3}
\input{sections/8_evaluation}

% === 参考文献 ===
\begin{thebibliography}{99}
% \bibitem{ref1} 作者. 标题[J]. 期刊, 年份.
\end{thebibliography}

% === 附录 ===
\begin{appendices}
\input{sections/A_code}
\end{appendices}

\end{document}
